% =====================================================================
%  CHAPTER 2: 生长发育概述（对应大纲第一章：生长发育规律部分 + 生长发育调查与评价）
%  参考PPT：第二章 儿童少年生长发育概述（苏普玉，安徽医科大学）
% =====================================================================
\chapter{生长发育概述}
\setcounter{chapter}{2}
\chapterintro{
掌握生长发育基本概念（生长、发育、成熟、生长发育可塑性）；掌握生长发育指标体系（体格发育指标、体能发育指标、心理行为发育指标）；掌握生长发育一般规律——阶段性与连续性统一（生长轨迹现象、追赶性生长）、程序性与时间性协调（头尾发展律、近侧发展律、个体差异性）、不同步性与多样性平衡（Scammon生长模式的四种类型及子宫生长模式）、生长发育的高度可塑性；掌握生长发育调查方法（横断面设计、监测设计、纵向设计、序列设计）及调查内容（人体测量、体能测试、心理发育调查）；掌握生长发育评价方法——发育水平评价（离差法、百分位数法、Z分法）、发育速度评价（年增长值、年增长率、生长监测图）、发育匀称度评价（体型匀称指标、身材匀称）、发育年龄评价（形态年龄、第二性征年龄、齿龄、骨龄）、营养状况评价（身高别体重、皮褶厚度）；熟悉生长发育调查应用、质量控制与伦理学；熟悉生长发育基本理论观点（进化观、社会生态学理论、整体观）。
}

% =====================================================================
% 第一部分：生长发育概念与指标体系
% =====================================================================
\section{生长发育概念与指标体系}

\subsection{生长发育基本概念}

\begin{keybox}
\textbf{（一）生长（Growth）}

身体各部分及全身在大小、长短和重量上的增加及身体化学成分的变化，含形态生长和化学生长。

\textbf{（二）发育（Development）}

身体组织、器官、系统在功能上不断分化与完善过程，含"身"（体格、体力）和"心"（心理、情绪、行为）两个密不可分的方面。

\textbf{（三）成熟（Maturity）}

生长和发育达到一个相对完备的阶段，标志个体在形态、生理功能、心理素质等方面都已达到成人水平，具备独立生活和生殖养育下一代的能力。

\textbf{（四）生长发育可塑性（Plasticity of Growth and Development）}

人的结构、功能为适应环境（积极/消极的内外环境）变化和生活经历发生改变的能力。
\end{keybox}

\subsection{生长发育指标体系}

\begin{defbox}
\textbf{（一）体格发育（Physical Growth）}

身体外部形态的发育，是人体整体发育的重要方面。

\begin{enumerate}[leftmargin=*]
    \item \imp{纵向测量指标}
    \item \imp{横向测量指标}：围度和径长。
    \item \imp{重量测量指标}：体重。
    \item \imp{体格发育派生指标}：BMI、腰高比、腰臀比。
\end{enumerate}
\end{defbox}

\begin{keybox}
\textbf{（二）体能发育指标}

\textbf{1. 生理功能指标：}
\begin{itemize}[leftmargin=*]
    \item \imp{心血管功能指标}：心率、脉搏、动态血压
    \item \imp{肺功能指标}：呼吸频率、肺活量、最大通气量、最大摄氧量
    \item \imp{肌肉发育指标}
\end{itemize}

\textbf{2. 运动能力指标：}
\begin{itemize}[leftmargin=*]
    \item \imp{力量指标}
    \item \imp{耐力指标}
    \item \imp{速度指标}
    \item \imp{灵敏性指标}
    \item \imp{柔韧性指标}
\end{itemize}

\textbf{3. 体能发育派生指标}
\end{keybox}

\begin{keybox}
\textbf{（三）心理行为发育指标}

\begin{itemize}[leftmargin=*]
    \item \imp{认知能力指标}：感知能力、记忆能力、注意能力、思维能力、执行能力
    \item \imp{情绪状态指标}：焦虑、抑郁、恐惧、偏执
    \item \imp{个体发育指标}
    \item \imp{社会适应能力指标}
\end{itemize}
\end{keybox}

% =====================================================================
% 第二部分：生长发育一般规律（大纲双横线——掌握）
% =====================================================================
\section{儿童少年生长发育一般规律}

\subsection{生长发育阶段性与连续性的统一}

\begin{keybox}
\textbf{（一）生长发育的阶段性}

\textbf{（二）生长发育的连续性}

\textbf{1. 生长轨迹现象（Trajectory of Growth）/ 生长管道现象（Growth Canalization）}

群体儿童少年在正常环境下，生长过程将按遗传潜能决定的方向、速度和目标发育。个体儿童的生长轨迹既与遗传有关，还与疾病及其治疗、营养、体力活动、情感环境等多种因素有关。

\textbf{2. 追赶性生长（Catch-up Growth）}

处于生长发育过程中的个体受疾病、营养、心理应激等因素作用下，会出现暂时性生长发育的连续性被打破现象，如生长迟缓，一旦这些影响因素解除，机体表现为向原有正常轨道靠近并具有生长发育强烈的倾向，年龄越小越明显，生长加速幅度越大。这种在阻碍生长发育因素解除后出现的生长加速现象。
\end{keybox}

\subsection{生长发育程序性和时间性的协调}

\begin{keybox}
\textbf{（一）生长发育的程序性}

\textbf{1. 运动能力发育的程序性}

\begin{itemize}[leftmargin=*]
    \item \imp{头尾发展律（Principal of Cephalocaudal Development）}：婴幼儿时期，粗大动作按抬头、翻身、坐、爬、站、走、跑、跳的发育程序进行。
    \item \imp{近侧发展律（Principal of Proximodistal Development）}：粗大动作和精细动作遵循近躯干的四肢肌肉先发育，手的精细动作后发育。
\end{itemize}

\textbf{2. 线性生长规律}

青春期前的身体线性生长也遵循头尾发展律，2个月龄的胎儿头与躯干比例1:1，出生时1:4，成人1:8。

\textbf{（二）生长发育的个体差异性}
\end{keybox}

\subsection{生长发育不同步性与多样性的平衡}

\begin{keybox}
\textbf{（一）不同组织器官系统的生长不同步性——斯卡蒙生长模式}

\begin{table}[H]
\centering
\caption{斯卡蒙（Scammon）生长模式的四种类型}
\begin{tabularx}{\textwidth}{lX}
\hline
\textbf{类型} & \textbf{特点} \\
\hline
\imp{一般型} & 婴幼儿期增长快；学龄前、学龄期增长平稳；青春期增长又加快；青春后期生长逐步停止 \\
\imp{淋巴系统型} & 淋巴结、胸腺、扁桃体等淋巴器官儿童期生长很快，青春期达高峰，以后逐渐衰退，成年时仅相当于高峰时的一半 \\
\imp{神经系统型} & 中枢神经系统，视听觉器官、颅骨等仅有一个生长突增期：出生前至学龄前期生长迅速，6岁左右成熟度达90\%左右 \\
\imp{生殖系统型} & 除子宫外的生殖器官，青春期前几乎呈停滞状态，青春期一开始迅速加快 \\
\hline
\end{tabularx}
\end{table}

\textbf{子宫生长模式（附加模式）：}子宫、肾上腺出生时较大，其后迅速变小，青春期开始前恢复到出生时的大小，其后迅速增大。

\textbf{（二）生长发育多样性}
\end{keybox}

\subsection{生长发育的高度可塑性}

\begin{defbox}
生长发育的高度可塑性：与年龄、环境和干预敏感期关系密切。
\end{defbox}

\subsection{生长发育规律的卫生学意义}

\begin{keybox}
\textbf{生长发育规律的卫生学意义：}

\begin{enumerate}[leftmargin=*]
    \item \imp{为制定卫生标准提供科学依据}：生长发育规律是制定学校课桌椅标准、教学作息制度、体育锻炼负荷等卫生标准的基础。了解各年龄段生长发育水平和速度，有助于科学确定卫生要求的适宜界值。

    \item \imp{指导生长发育监测与评价}：掌握生长轨迹现象和追赶性生长规律，有助于正确判断生长发育偏离，避免将暂时性生长迟缓误判为永久性障碍，指导早期干预的时机和策略。

    \item \imp{指导学校健康教育}：根据生长发育的程序性和阶段性特点，在不同年龄阶段开展有针对性的健康教育和健康促进活动，使教育内容与学生的发育水平相适应。

    \item \imp{为疾病预防提供依据}：了解不同组织器官系统的生长模式（Scammon生长模式），有助于识别各系统疾病的高发年龄段，制定针对性的预防策略。如淋巴系统型提示青春期前是免疫相关疾病防控的关键期。

    \item \imp{促进健康政策制定}：生长发育的长期变化趋势（secular trend）是评价社会经济发展、营养改善和公共卫生干预效果的敏感指标，为制定国家层面的儿童健康政策提供决策依据。

    \item \imp{体现生命历程健康观}：生长发育的高度可塑性表明，生命早期的良好环境和干预可对终身健康产生深远影响，体现了DOHaD理论和全生命周期保健策略的重要性。
\end{enumerate}
\end{keybox}

% =====================================================================
% 第三部分：生长发育调查
% =====================================================================
\section{生长发育调查}

\subsection{生长发育调查概述}

\begin{defbox}
\textbf{生长发育调查：}运用科学设计和方法对个体或群体儿童少年的生长发育状况进行观察和测量，以研究生长发育的规律和影响因素。

\textbf{目的：}
\begin{enumerate}[leftmargin=*]
    \item 研究不同儿童少年群体的生长发育规律和特点。
    \item 制定国家、地区儿童少年生长发育正常值或评价标准。
    \item 探索各种内外因素对生长发育的影响，发现主要影响因素及机制，提出相应干预措施。
    \item 检验和评价各种学校保健、干预措施对生长发育的保障或促进作用。
    \item 比较不同人群、不同历史时期的生长发育资料，研究变化趋势和影响因素。
\end{enumerate}
\end{defbox}

\subsection{生长发育调查设计}

\begin{keybox}
\textbf{（一）横断面设计（Cross-sectional Design）}

在某一时间段内，选择特定地区、有代表性的对象即不同年龄阶段的儿童少年，针对特定指标，进行一次性的群体大规模的调查。是生长发育调查的主要方法，如全国学生体质健康调研。

\begin{itemize}[leftmargin=*]
    \item 可在短期内获得大量数据。
    \item 应有周详的调查方案，选取的指标要少而精。
    \item 应采取适宜的抽样方案，如多阶段、分层随机整群抽样，使对象有较强代表性。
    \item 参与测试人员较多，检测方法和技术规范应统一。
\end{itemize}

\textbf{（二）监测设计（Surveillance Design）}

对某地区、某群体的生长发育指标的连续收集、整理、分析过程。

\begin{enumerate}[leftmargin=*]
    \item 明确监测目的。
    \item 确定监测人群：可在学校、社区、医院等不同基础上对学龄儿童少年进行人群监测（population-based surveillance）。
    \item 稳定监测方法；监测体系固定，指标简便易行。
    \item 保证质量控制；保障结果的可比性。
    \item 结果反馈与应用；及时制定和调整干预措施。
\end{enumerate}

\textbf{（三）纵向设计（Longitudinal Design）/ 随访设计（Follow-up Design）}

在一段时间如几年乃至终生对同一批儿童少年个体进行反复观察的研究设计。

\begin{itemize}[leftmargin=*]
    \item \imp{优点}：研究对象是同一组个体，因而能揭示早期经历（不良环境、营养不良等）与后来生长发育结果间的因果关系。
    \item \imp{缺点}：耗时长，花费大，样本易流失。
\end{itemize}

\textbf{（四）序列设计（Sequence Design）}

对不同年龄组的儿童少年进行横向调查，间隔一定时间后，对同一批儿童少年进行一次或多次重复调查，构成追踪性研究，将横断面调查和追踪调查结合。

\begin{itemize}[leftmargin=*]
    \item \imp{优点}：①克服了追踪调查所需年限太长、样本易流失等缺点；②节约时间和工作量。
    \item \imp{缺点}：仅具部分追踪性质，所获生长速度数据是近似的。
\end{itemize}
\end{keybox}

\subsection{生长发育调查内容}

\begin{defbox}
\textbf{（一）身体发育调查}

身体发育调查包括：体格指标、生理功能指标、运动素质指标、皮褶厚度、身体成分、体型、骨发育、性发育。

\textbf{1. 人体测量：}对人体有关部位长度、宽度、厚度和围度的测量，是评价儿童身体匀称度、大小和成分的最简便指标，反映儿童健康和营养状况。

\begin{itemize}[leftmargin=*]
    \item \imp{长度}：身高、体重（最常用）、坐高、手长、足长、上肢长、下肢长。
    \item \imp{宽度}：肩宽、骨盆宽、胸廓横径、胸廓前后径、指间距。
    \item \imp{围度}：头围、胸围、腰围、腹围、上臂围、大腿围、小腿围。
\end{itemize}

\textbf{2. 体能测试：}采用医学手段检测机体在日常体力活动或运动中的功能。

\begin{itemize}[leftmargin=*]
    \item \imp{身体成分}：人体中脂肪、骨组织、肌肉组织占人体总质量的百分比，如体质指数、皮褶厚度、生物电阻抗分析。
    \item \imp{心肺耐力}：反映心血管系统和呼吸系统在长时间体力活动中耗氧情况和长时间剧烈运动能力，反映体能与健康关系、直接指示人的生理状况，如肺活量、心率、脉率、血压、20m跑、3min阶梯测试。
    \item \imp{肌力肌耐力}：握力、背肌力、立定跳远。
    \item \imp{柔韧性}：屈体前伸测试、肩力量测试。
\end{itemize}

△肺活量指数：男青春期前↑后不变，女青春期前↑后↓。

△素质发育顺序：速度→耐力→腰腹→下肢→上臂。

\textbf{（二）心理发育调查}

心理发育包括认知、情绪、个性和行为等。

\textbf{1. 智力测验：}通过心理测验量表获得有关儿童智力水平和能力特点等信息，为诊断儿童智力发育异常提供依据。

\begin{itemize}[leftmargin=*]
    \item \imp{筛查性测验}
    \begin{itemize}
        \item 项目：丹佛发育筛查测验、绘人试验、图片词汇测验和学前儿童能力筛查（"50项"）。
        \item 特点：简便快速，可大致筛出正常与异常（可疑）儿童，但不能判断儿童异常程度。
    \end{itemize}
    \item \imp{诊断性测验}
    \begin{itemize}
        \item 项目：贝利婴幼儿发育量表、Gesell发育量表、斯坦福-比奈量表、韦氏学前儿童智力量表、韦氏儿童智力量表。
        \item 特点：内容全面复杂，测验结果能较精确客观反映受试者心理行为发育水平，可作为评价儿童发育水平的重要依据之一。
    \end{itemize}
\end{itemize}

\textbf{2. 行为评定：}通过观察、检查、量表等方法对个体行为问题作出评价，以尽早开展问题矫治，对保证儿童身心健康发展有重要意义。

\begin{itemize}[leftmargin=*]
    \item 项目：标准化评定量表。
    \item 评价方法：他评或自评。
    \item 多数量表可评定儿童多项行为或行为问题，少数量表只评定儿童单项行为。
\end{itemize}
\end{defbox}

\subsection{生长发育调查应用}

\begin{defbox}
\begin{enumerate}[leftmargin=*]
    \item \imp{描述水平}
    \item \imp{反映动态}：对同地区、同人群进行连续多次调查可追踪儿童少年生长发育动态变化，分析生长长期趋势。
    \item \imp{评价队列效应}
\end{enumerate}

△\textbf{队列效应（Cohort Effect）/ 同辈效应：}某群体与其他不同年龄段群体间的差异是由这个群体成长时相似的文化背景、社会风俗、价值观、经历、受教育程度、生活习惯等影响造成的，而不是由真正的心智上的发展引起的。即不同出生队列因为经历不同时期，其生理及心理健康发展状况可反映相应阶段的社会变化与时代特征。

\textbf{4. 制定参考标准}

\textbf{（1）建立正常值}

\begin{itemize}[leftmargin=*]
    \item 传统方法：均数±标准差。
    \begin{itemize}
        \item \imp{现状正常值}：样本来自本地区中等水平人群，只剔除那些因慢性疾患或残障而对生长发育有明显不良影响的个体。反映本地区、本人群的真实情况，也是干预目标。
        \item \imp{理想正常值}：对象都应是足月正常体重新生儿；自幼生活在适宜环境，营养充分，无慢性病史，有良好生活条件和卫生服务；生长水平高，生长潜力得以充分发挥。一旦建立，能长期使用。
    \end{itemize}
    \item 目前国际常用方法：采用经LMS法（实现资料正态化）修正后的百分位数法。按P5P25P50P75P95分5个等级；在此基础上增加P3P10P90P97形成9个等级。适用于筛查异常和"疑似异常"。
\end{itemize}

△我国目前难以像欧美国家那样找到"理想人群"，建立现状正常值更现实。

\textbf{（2）生长发育标准}

\begin{itemize}[leftmargin=*]
    \item 标准可看成正常值的一部分，制定样本应尽量接近"理想人群"。
    \item 全国性评价标准和地区性评价标准各有优势，应根据目的选择，配合使用。
\end{itemize}
\end{defbox}

\subsection{生长发育调查质量控制}

\begin{keybox}
\textbf{（一）调查实施细则}

\begin{enumerate}[leftmargin=*]
    \item \imp{检测仪器和方法}：①检测仪器应精确，检测方法应统一；②追踪调查应自始至终使用同一方法和同种仪器。
    \item \imp{检测时间}：①追踪调查时个体测量时间应相对固定；②将同龄样本均匀分配在上下午；③考虑季节、生活制度对生长发育的影响，一般5-6月或9-10月。
    \item \imp{年龄}：①我国按实足年龄（测试年月日-出生年月日）计算，满7岁-差1天满8岁者都为7岁；②不同国家算法不同，比较国别资料时应注意消除该差异。
    \item \imp{技术培训}：严格培训测试者，考核合格方可上岗。
    \item \imp{检测程序}：各项目按规定顺序流水作业，以免漏测。
\end{enumerate}

\textbf{（二）质量控制}

\begin{enumerate}[leftmargin=*]
    \item 人员检验
    \item 现场检验：①逐一核对调查表，检查缺项、漏项；②记录准确，字迹清晰；③计算误差率：A.若P>5\%，应及时分析原因并采取措施改进；B.若P>20\%，当日全部检测数据无效，重新测试。
    \item 逻辑检验
\end{enumerate}
\end{keybox}

\subsection{生长发育调查伦理学}

\begin{defbox}
\begin{enumerate}[leftmargin=*]
    \item \imp{伦理审批}：向当地/高校伦理委员会提交申请报告。
    \item \imp{知情同意}：以口头/书面形式，取得对象（14岁以下儿童征询家长意见）的参与承诺。
    \item \imp{儿童少年的禁忌}：凡对实验对象可能造成伤害的试剂或药品（如放射性核素）不得用于儿童少年。
    \item \imp{正确对待对照组}。
\end{enumerate}
\end{defbox}

% =====================================================================
% 第四部分：生长发育评价
% =====================================================================
\section{生长发育评价}

\subsection{生长发育评价概述}

\begin{keybox}
\textbf{生长发育评价（Appraisal of Growth and Development）：}把儿童少年各项生长指标的实测值与参考标准进行比较，以分析儿童少年生长发育水平、变化和个体差异，分析发育的趋势等。

\textbf{意义：}
\begin{enumerate}[leftmargin=*]
    \item 有助于了解和评价个体或群体儿童少年现时生长发育水平、今后发展趋势等。
    \item 筛查、诊断生长发育障碍，评价营养和生活环境因素对生长发育的影响，提供保健咨询建议。
    \item 列入社区健康水平指标系，通过观察指标变化，评价各项学校卫生措施的实效，作为实施学校卫生监督的依据。
    \item 从群体角度比较不同人群、地区及年代差异，分析社会经济发展对儿童少年生长发育的影响。
\end{enumerate}
\end{keybox}

\subsection{生长发育评价的内容与参考标准}

\begin{defbox}
\textbf{内容：}体格生长评价、体能发育评价、性发育评价、行为发育评价、智力发育评价。

\textbf{目标：}
\begin{enumerate}[leftmargin=*]
    \item 生长水平
    \item 生长速度
    \item 匀称度
\end{enumerate}

\textbf{生长发育参考标准：}
\begin{enumerate}[leftmargin=*]
    \item 5岁↓儿童建议采用WHO生长标准。
    \item 学龄期，尤其是青春期开始后出现较大差异，中国相对略高的界值点百分分位线可能利于早期筛查生长异常的儿童。
\end{enumerate}
\end{defbox}

\subsection{生长发育水平评价}

\begin{keybox}
\textbf{（一）离差法}

\textbf{1. 评价原理：}用均值($\bar{x}$)和标准差($s$)描述样本调查值的分布，适用于正态分布资料。以均值($\bar{x}$)±标准差($s$)表示。$\bar{x}\pm s$包括样本的68.3\%，$\bar{x}\pm 2s$包括样本的95.4\%，$\bar{x}\pm 3s$包括样本的99.7\%。一般以$\bar{x}\pm 2s$为正常范围。

\textbf{2. 评价方法}

\textbf{（1）等级评价法：}用标准差与均值的位置远近划分等级，国内常用五等级评价标准。

\begin{itemize}[leftmargin=*]
    \item \imp{应用}
    \begin{itemize}
        \item 评价个体：将个体发育指标的实测值与同年龄、同性别正常值比较，确定发育等级。
        \item 评价群体：先将2个班/学校所有个体按标准逐个划入不同等级，再计算各等级人数及构成比。
    \end{itemize}
    \item \imp{注意事项}
    \begin{itemize}
        \item 正确理解等级划分的含义，不是越高越好。
        \item 个体实测值在$\bar{x}\pm 2s$内，视为正常；在$\bar{x}\pm 2s$外，不能轻易判定为异常，需连续观察、结合临床检查，慎重下结论。
    \end{itemize}
    \item \imp{优点}：方便快捷，适用于基层。
    \item \imp{缺点}：①只适用于正态分布指标，对偏态分布指标（如体重、BMI）评价不够精确。②结果只能针对单项指标，无法反映个体发育匀称程度。③不能直观反映动态变化。
\end{itemize}

\textbf{（2）曲线图法}

\begin{itemize}[leftmargin=*]
    \item 应用：通过对同一年龄组儿童少年不同年代发育均值比较，可分析长期趋势。
    \item 优点：使用简便、结果直观，可针对个体和群体。
    \item 缺点：需要分男女，事先准备好针对不同指标的分年龄图。
\end{itemize}

\textbf{（二）百分位数法}

\textbf{1. 评价原理：}按某一指标（如身高、体重）不同个体测量值按从小到大顺序排列，分为100等份，每1等份即代表1个百分位值。适用于正态或非正态分布调查资料，一般采用P3P10P25P50P75P90P97为主要界值点。

\textbf{2. 评价方法}

\textbf{（1）等级评价法}

\begin{table}[H]
\centering
\caption{离差法与百分位数法等级对应关系}
\begin{tabularx}{\textwidth}{lll}
\hline
\textbf{等级} & \textbf{离差法} & \textbf{百分位数法} \\
\hline
上等   & $>\bar{x}+2s$     & $>P_{97}$ \\
中上等 & $>\bar{x}+s\sim\bar{x}+2s$ & $>P_{75}$ \\
中等   & $\bar{x}\pm s$           & $P_{50}$ \\
中下等 & $<\bar{x}-s\sim\bar{x}-2s$ & $<P_{25}$ \\
下等   & $<\bar{x}-2s$     & $<P_3$ \\
\hline
\end{tabularx}
\end{table}

\textbf{（2）曲线图法}

\begin{itemize}[leftmargin=*]
    \item 应用：评价群体儿童时，可单用某指标的P50或配合P10P25P75P90等少量曲线，比较同时期不同群体发育水平差异或比较同群体不同年代变化趋势。
    \item 优点：形象直观，反映发育水平准确，便于动态观察。
\end{itemize}

\textbf{（三）Z分法（未实现资料正态化）}

\textbf{1. 评价原理：}一种特殊类型的离差法，以0为中心，用偏离该年龄组标准差的程度反映生长情况。

\textbf{2. 评价方法：}$Z=(x-\bar{x})/s$

\begin{table}[H]
\centering
\caption{Z分法等级划分}
\begin{tabularx}{\textwidth}{llllll}
\hline
\textbf{等级} & 上等 & 中上等 & 中等 & 中下等 & 下等 \\
\hline
\textbf{评分} & $>2$ & $1\sim2$ & $-1\sim1$ & $-2\sim-1$ & $<-2$ \\
\hline
\end{tabularx}
\end{table}
\end{keybox}

\subsection{生长发育速度评价}

\begin{keybox}
\textbf{（一）群体生长速度评价}

\textbf{1. 年增长值：}群体两连续年龄组身高均值相减。

\textbf{2. 年增长率：}年增长值/身高基数。

\begin{itemize}[leftmargin=*]
    \item 年增长率：$V(\%)=(H_{t+1}-H_t)/H_t\times 100\%$。
    \item 修匀值：$V=[(a+b)/2+(b+c)/2]/2=(a+2b+c)/4$，$b$：修正后该年龄组年增长率，$a$、$c$：$b$上下相邻年龄组年增长率。
\end{itemize}

\textbf{（二）个体生长速度评价}

常用生长监测图表示，所用正常值标准应以追踪性资料为基础建立。结果以正常、下降、缓慢（增长不足）、加速表示。
\end{keybox}

\subsection{发育匀称度评价}

\begin{keybox}
\textbf{（一）体型匀称度}

\textbf{1. 身高体重指数/克托莱指数} $=$ 体重(kg)/身高(cm)$\times$100\%。

意义：单位身高的体重，反映人体充实度和营养状况，随年龄↑而↑，男性21岁、女性19岁趋于稳定。

\textbf{2. 身高胸围指数} $=$ 胸围(cm)/身高(cm)$\times$100\%。

意义：反映胸廓发育状况和从横截面反映躯干体型。均值在突增高峰前随年龄↑而↓，突增高峰时最低；突增高峰后随年龄↑而↑，成年后趋于稳定。

\textbf{3. 体质量指数（BMI）} $=$ 体重(kg)/[身高(m)]$^2$。

意义：每平方米身体面积包含的体重，敏感反映人体充实度和体型胖瘦，受身高干扰小，与皮脂厚度、上臂围等营养指标相关性较高，广泛用于建立营养不良、超重/肥胖筛查指标，成人$\geq$24超重、$\geq$28肥胖。

△\textbf{Kaup指数} $=$ 体重(kg)/[身长(cm)]$^2\times 10^4$。意义：反映婴幼儿体格发育和营养状况、儿童机体脂肪总量及分布情况的指标。

\textbf{4. 劳雷尔指数（Rohrer指数）} $=$ 体重(kg)/[身高(cm)]$^3\times 10^7$。

意义：每cm$^3$身体的重量。曲线呈"V"字形，5岁后随年龄↑而↑，能较敏感反映体型胖瘦。缺点是受身材高矮影响大，如男、女12岁左右有一次交叉，交叉前男>女，交叉后女>男；据此判断营养状况有明显不合理处。

\textbf{（二）身材匀称}

\begin{itemize}[leftmargin=*]
    \item 以坐高（3岁前儿童顶臀高）/身高（3岁前儿童身长）比值（SH/H）或躯干/下肢反映下肢发育状况。
    \item 从婴儿期的0.68逐渐下降至青少年的0.52，提示青春期前下肢较躯干生长快，SH/H与身高呈显著负相关。
    \item 用坐高、身高测量值计算比值与参考人群比值比较，实际比值$\leq$参照人群值为身材匀称，实际比值$>$参照人群值为不匀称。
    \item 身材匀称的评价结果可帮助诊断内分泌及骨骼发育异常疾病。
\end{itemize}

△等级评价法：发育水平；曲线图法：发育水平、速度趋势；相关回归法：发育水平、匀称程度。
\end{keybox}

\subsection{发育年龄评价}

\begin{keybox}
\textbf{（一）形态年龄（Morphological Age）}

用一般群体某形态指标发育状态制成标准年龄，将个体该形态指标与标准指标比较来评价，最常用身高年龄和体重年龄。

\begin{itemize}[leftmargin=*]
    \item \imp{优点}：使用简单，结果明确。
    \item \imp{缺点}：①单项形态指标仅反映全身发育一个侧面；②单独评价效果常不理想，仅在衡量该体格指标的发育提前或滞后时有提示作用。
\end{itemize}

\textbf{（二）第二性征年龄（Secondary Sex Characteristic Age）}

用第二性征发育指标制成的标准年龄，根据个体第二性征发育程度与其比较来评价。

\begin{itemize}[leftmargin=*]
    \item 男性：睾丸容积、阴毛、喉结、变声；女性：乳房、腋毛和阴毛。
    \item 第二性征年龄评价只适用于青春期（初高中）。
\end{itemize}

\textbf{（三）齿龄/牙齿年龄（Dental Age）}

按儿童牙齿发育顺序制成标准年龄，反映个体发育状况（小学）。评价方法：

\begin{enumerate}[leftmargin=*]
    \item \imp{牙齿萌出数量和质量}：适用于出生后6个月-13岁左右（除第三磨牙外），对早期发育评价有重要价值。
    \item \imp{X线摄片法}：第1颗牙钙化开始-最后1颗恒牙完成钙化的全过程，划分不同阶段反映牙齿发育水平和速度。该方法准确可靠。
\end{enumerate}

\textbf{（四）骨龄/骨骼年龄（Skeletal Age）}

骺和小骨骨化中心出现的年龄及骺及骨干愈合的年龄，是反映个体发育水平和成熟程度的精确指标，根据儿童少年骨骼钙化程度与骨发育标准比较可评价其发育年龄。以手腕部最理想，优点：

\begin{enumerate}[leftmargin=*]
    \item 手、腕骨数目、种类和形状多样，包括长骨、短骨、不规则骨和籽骨，对全身骨骼有很好代表性。
    \item 手、腕骨各继发性骨化中心出现及掌指骨、尺桡骨的干骺愈合有明显时间顺序，不同发育阶段间界限明确，易发现差别。
    \item 拍片方便，投照条件易控制，受检者接受X线剂量小，对保护儿童少年健康有利。
\end{enumerate}
\end{keybox}

\subsection{营养状况评价法}

\begin{keybox}
\textbf{（一）身高别体重}

\begin{itemize}[leftmargin=*]
    \item WHO推荐的指标，反映儿童现时营养状况，筛查营养不良。
    \item 在同等身高条件下比较体重大小，可有效消除青春期前因性别、发育水平、遗传、种族差异等原因导致的身材发育差异的影响。
    \item \imp{优点}：使用简便，评价的营养水平较准确、灵敏和客观。
    \item 目前在各发达国家，针对学龄儿童少年群体超重肥胖筛查，身高别体重已逐步被BMI（体重指数）标准取代。
\end{itemize}

\textbf{（二）皮褶厚度}

\begin{itemize}[leftmargin=*]
    \item 通过估测皮下脂肪（约占全身脂肪量50\%↑）反映儿童少年近期营养状况的方法之一，用于评价肥胖程度效果较好。
    \item 测量部位有多处，其中以上臂肱三头肌部（代表四肢）和肩胛下角部（代表躯干）最理想，两部位测量值之和可代表全身皮下脂肪发育状况。
    \item 皮褶厚度和体脂含量间相关性较高，回归系数在0.7左右，可用皮褶厚度建立估计体脂含量百分比（即脂肪含量占体重的\%，简称体脂率）的回归方程。
\end{itemize}
\end{keybox}

% =====================================================================
% 第五部分：生长发育基本理论观点
% =====================================================================
\section{生长发育基本理论观点}

\subsection{生长发育进化观}

\begin{defbox}
\textbf{1. 进化形成的适应机制在生长发育中的作用}

\textbf{2. 生命史策略在生长发育中的作用：}生命史理论研究生物体在不同的生命目标（如生存、生长和繁殖）之间进行权衡并适应性地分配资源的理论。个体在资源有限的情况下，需要考虑如何分配自身资源，而个体早期成长环境会通过内分泌轴，影响资源在生存、生长和繁衍的分配（即生命史策略）。
\end{defbox}

\subsection{社会生态学理论}

\begin{keybox}
\textbf{Bronfenbrenner生态系统理论：}

生态系统理论是发展心理学的重要理论之一，由美国心理学家尤里·布朗芬布伦纳（Urie Bronfenbrenner，1917—2005）建立的个体发展模型。他认为个体的发展嵌套于相互影响的一系列的环境系统中，各个系统之间存在交互耦合的联系和影响，个体与环境相互作用并影响着个体的发展。

\textbf{（一）空间维度：}
\begin{enumerate}[leftmargin=*]
    \item \imp{微系统（Microsystem）}：是个体活动和人际交往的直接环境，不断在变化和发展，包括家庭、学校、同伴群体、社区等。
    \item \imp{中系统（Mesosystem）}：指各微系统间的联系或相互关系，如家庭、同伴和学校之间的关系。
    \item \imp{外系统（Exosystem）}：指儿童并未直接参与但却对他们的发展产生影响的环境。
    \item \imp{宏系统（Macrosystem）}：指存在于以上三个系统中的文化、亚文化和社会环境。
\end{enumerate}

\textbf{（二）时间维度（历时系统，Chronosystem）：}

布朗芬布伦纳还把时间维度纳入模型，将时间作为研究个体成长中的心理变化参照体系，称为历时系统。他强调儿童的发展、变化是一个将时间和环境结合起来考察的动态过程。生态系统不是一种恒定不变的系统，它始终处于动态变化中。
\end{keybox}

\subsection{整体观}

\begin{defbox}
\textbf{生长发育整体观}主张将儿童少年的生长发育人为地划分成几个学科领域，不同学科领域在揭示生长发育现象时，本质相同但各有侧重点。生长发育整体观理论的核心观念是：应将儿童少年生长发育的各个方面都视为是存在有机联系的，相互依赖的，而不是零碎的、割裂的。每个人在生物性、认知性和社会性等方面都具有自身的特征，每个人又受制于各个层次系统"大环境"的直接间接的影响。
\end{defbox}

\newpage
